\chapter{Coexistence and Gradual Migration Strategy}

The transition from the legacy SRP system to the new ServiceRequestManager architecture represents a substantial technical and operational challenge that requires careful planning to minimize business disruption while ensuring data integrity and system reliability. Rather than implementing a disruptive one-shot migration, the project adopts a phased coexistence strategy that enables gradual migration to the new system while maintaining full operational continuity.

\section{Migration Strategy Overview}

The migration architecture implements a dual-flow execution model where both SRP and SRM systems operate simultaneously, with individual service request actions configured to route through either the legacy or modern execution pipeline based on the \texttt{target\_type} enumerated attribute of the Action model.

% \begin{listing}[htbp]
% \begin{minted}[
%   linenos,
%   tabsize=1,
%   fontsize=\footnotesize,
%   breaklines
% ]{python}
% class ActionTargetTypes(Enum):
%     DOMAIN_DC = "From SR Domain (Domain Controller)"
%     DOMAIN_EXC = "From SR Domain (Exchange Server)"
%     INTERNAL_WORKER = "Internal Worker"
%     INTERNAL_API = "API Call"
%     SRP = "SRP Agent (Domain Controller)"  # Legacy flow
% \end{minted}
% \caption{Action Target Type Enumeration for Dual-Flow Routing}
% \label{lst:target_type_migration}
% \end{listing}

Actions configured with \texttt{target\_type = SRP} continue routing through the legacy SRP Agent execution mechanism, preserving the existing workflow without modification. Actions configured with any other target type utilize the new Ansible-based execution pipeline through Provision Prime, benefiting from improved logging, error handling, and state management capabilities introduced in SRM.

This per-action migration control provides several strategic advantages: it enables pilot deployments on specific customers or action types; it maintains business continuity by ensuring that any discovered issues affect only the subset of actions migrated to the new flow; it supports progressive learning; and it provides fallback capability, as actions can be reverted to SRP execution if critical issues are discovered.

The primary architectural trade-off is that the catalog configuration must be maintained in the SRP database during the coexistence period, as SRP lacks the sophisticated configuration features introduced in SRM (three-level hierarchy, external API configurations, template-based autofill). The SRM system operates as a read-only consumer of catalog data from SRP, applying its advanced configuration features as an overlay on the imported base catalog.

\section{Catalog Synchronization from SRP to SRM}

The catalog synchronization mechanism implements a unidirectional data flow where SRP serves as the authoritative source of catalog definitions, with SRM periodically importing and enriching this data. This synchronization is executed through the \texttt{sync\_srp\_catalog} Django management command, scheduled as a Kubernetes CronJob for daily execution at 01:00 AM.

% \begin{listing}[htbp]
% \begin{minted}[
%   linenos,
%   tabsize=1,
%   fontsize=\footnotesize,
%   breaklines
% ]{python}
% @monitor_job
% def handle(self, *args, **kwargs):
%     with disable_auditlog(), pyodbc.connect(SRP_CONN_STR) as conn:
%         cursor = conn.cursor()
%         sync_action_types(cursor)
%         sync_action_categories(cursor)
%         sync_action_groups(cursor)
%         sync_customers(cursor)
%         sync_actions(cursor)
%         sync_config_parameters(cursor)
%         sync_fields(cursor)
%         # ...
%         reset_sequences(app_label="catalog")
% \end{minted}
% \caption{Catalog Synchronization Job Structure}
% \label{lst:sync_catalog_job}
% \end{listing}

\subsection{Synchronization Architecture}

The synchronization process establishes a direct ODBC connection to the SRP SQL Server database and executes a series of import operations that reconstruct the complete catalog hierarchy in the SRM PostgreSQL database. The process imports entities in a carefully ordered sequence that respects foreign key dependencies: action types → categories → groups → customers → actions → configuration parameters → fields → customer-specific customizations.

The synchronization executes with audit logging disabled to prevent creation of thousands of audit log entries for bulk import operations, improving performance and reducing database bloat. Each synchronization function implements a consistent pattern: query the source table in SRP, construct Django model instances for new or modified records, and execute bulk create/update operations.

\subsection{Differential Synchronization Logic}

The synchronization functions implement a differential logic that minimizes database operations by detecting actual changes rather than blindly recreating all data. Listing \ref{lst:differential_sync} demonstrates an example of the pattern used across all synchronization functions.

\begin{listing}[htbp]
\begin{minted}[
  linenos,
  tabsize=1,
  fontsize=\footnotesize,
  breaklines
]{python}
records = [dict(zip(columns, row)) for row in cursor.fetchall()]
source_ids = {int(r["IDCategory"]) for r in records}

existing_categories = {
    ac.id: ac for ac in ActionCategory.objects.filter(
        id__in=source_ids
    )
}

to_create = []
to_update = []

for record in records:
    category_id = int(record["IDCategory"])
    if category_id in existing_categories:
        obj = existing_categories[category_id]
        obj.description = record["DescCategory"]
        # ...
        to_update.append(obj)
    else:
        to_create.append(ActionCategory(
            id=category_id,
            description=record["DescCategory"]
            # ...
        ))

if to_create:
    ActionCategory.objects.bulk_create(to_create)
if to_update:
    ActionCategory.objects.bulk_update(to_update, 
        ['description', 'display_order'])
\end{minted}
\caption{Differential Synchronization Pattern with Bulk Operations}
\label{lst:differential_sync}
\end{listing}

This approach constructs in-memory dictionaries of existing records keyed by primary key, then iterates through source records to segregate them into create and update lists. Bulk operations execute with batch sizes appropriate for PostgreSQL, significantly outperforming individual save operations. An initial full synchronization strategy that deleted and recreated all records at each execution proved inefficient with SRP's~large data volumes, resulting in execution times around 2 hours. The current differential logic significantly reduces database load, bringing average execution times down to approximately 20 seconds.

\subsection{Post-Import Enrichment}

Following the base catalog import from SRP, the synchronization process executes several enrichment operations that apply SRM-specific configurations: API configuration synthesis creates \texttt{ExternalAPIConfiguration} records defining option sources for select fields; field duplication resolution identifies and resolves scenarios where multiple fields with identical names exist within the same action scope; HDA field configuration synthesizes metadata for external fields and classification fields; and customer-specific customizations translate SRP configurations into SRM's~three-level hierarchy.

The synchronization concludes by resetting PostgreSQL auto-increment sequences to prevent primary key conflicts on subsequent manual record creation.

\section{Execution History Synchronization}

While catalog synchronization flows unidirectionally from SRP to SRM, execution history must be merged from both systems to provide unified visibility. The \texttt{sync\_srp\_sr\_history} command implements this bidirectional synchronization, scheduled to execute at 02:00 AM (one hour after catalog sync).

The execution history synchronization imports three primary entity types: service request headers, field values, and execution logs. The synchronization is filtered by a configurable date threshold to avoid importing the complete historical archive, defaulting to the first day of the current month.

% \begin{listing}[htbp]
% \begin{minted}[
%   linenos,
%   tabsize=1,
%   fontsize=\footnotesize,
%   breaklines
% ]{python}
% start_date = os.getenv(
%     "SRP_SR_SYNC_START_DATE",
%     datetime.date.today().replace(day=1).strftime("%Y-%m-%d")
% )
% cursor.execute(
%     f"SELECT * FROM TBSR_Headers WHERE Date > '{start_date}'"
% )
% # ... process service requests
% sync_service_requests(cursor)
% sync_service_request_fields(cursor)
% sync_service_request_logs(cursor)
% \end{minted}
% \caption{Execution History Synchronization with Date Filtering}
% \label{lst:history_sync}
% \end{listing}

Service request field values imported from SRP must respect SRM's~enhanced security model where sensitive fields are encrypted at rest. The synchronization implements transparent encryption during import, checking each field's~sensitivity flag and applying encryption to values from sensitive fields. Execution logs stored in SRP use a numeric type code system incompatible with SRM's~enumerated log types, requiring a mapping table that translates SRP codes to SRM equivalents (Creation, Approval, Execution Started, Execution Completed, etc.).

\section{Dual Execution Flow Management}

The coexistence architecture's~core mechanism is the dynamic routing of service request execution based on action configuration. When a service request is created through the SRM API, the system evaluates the action's~\texttt{target\_type} to determine whether to invoke the new or legacy execution pipeline, as illustrated in Listing \ref{lst:dual_flow}.

\begin{listing}[htbp]
\begin{minted}[
  linenos,
  tabsize=1,
  fontsize=\footnotesize,
  breaklines
]{python}
is_srp_flow: bool = action.target_type == ActionTargetTypes.SRP.name

# Create HDA ticket (common to both flows)
response_ticket: dict = mysupport.create_ticket(
    codicli=codicli, idrs=customer.idrs, body=body
)
ticket_id: str = response_ticket.get("ID")

if is_srp_flow:
    # Legacy flow: SRP will create the SR entity
    logger.info(f"SRP flow detected. Ticket {ticket_id} created.")
    return ServiceRequest(
        action=action, customer=customer, 
        ticket_reference=ticket_id
    )
else:
    # Modern flow: SRM creates SR and triggers Celery tasks
    validated_data["ticket_reference"] = ticket_id
    service_request: ServiceRequest = super().create(validated_data)
    # ... trigger async execution
    return service_request
\end{minted}
\caption{Dual Execution Flow Routing Logic}
\label{lst:dual_flow}
\end{listing}

For actions configured with \texttt{target\_type = SRP}, the serializer creates the HDA ticket but returns a non-persisted ServiceRequest object. The actual service request record will be created by SRP's~legacy REST API when HDA invokes it, then synchronized into SRM's~database by the nightly job. For actions with any non-SRP target type, the ServiceRequest record is persisted immediately with NEW status, and the asynchronous Celery task chain is initiated.

The dual-flow architecture remains transparent to frontend applications, which interact exclusively with SRM APIs regardless of the underlying execution mechanism, ensuring that frontend applications require no modifications as actions are gradually migrated.

\section{Pilot Strategy and Progressive Rollout}

The migration strategy implements a risk-managed progression that validates the new architecture incrementally before full deployment, organized into four distinct phases.

\subsection{Phase 1: Internal Pilot}

The initial migration phase configures a small subset of actions used exclusively by internal IT operations teams to execute through the SRM pipeline. These actions are selected based on low execution volume, non-critical business processes, and comprehensive logging requirements. Internal teams provide direct feedback on execution behavior, enabling rapid iteration on issues discovered during real-world usage.

\subsection{Phase 2: Customer Pilot}

Following successful internal validation, the migration expands to a carefully selected cohort of pilot customers chosen based on technical sophistication, diverse action usage patterns, and strong relationship with the company. Each pilot customer has a subset of their actions migrated to SRM execution, typically starting with low-risk operations before progressing to more critical provisioning workflows.

The pilot phase implements enhanced monitoring including dedicated Grafana dashboards, daily review meetings, and direct communication channels with pilot customers. Issues discovered during pilot are addressed through hotfix deployments, with actions reverted to SRP execution if critical problems cannot be resolved quickly.

\subsection{Phase 3: Progressive Expansion}

Successful pilot completion triggers progressive expansion where additional customer cohorts are migrated on a weekly cadence. Each expansion wave includes a defined set of customers and actions, with execution metrics monitored for anomalies before proceeding to the next wave. The expansion velocity is dynamically adjusted based on observed success rates: high success rates (>99\%) enable acceleration, while elevated error rates trigger pause for investigation.

The rollout maintains a progressive migration of action types from simple to complex: read-only information queries migrate first, followed by non-destructive provisioning operations, then critical operations, and finally complex multi-step workflows.

\subsection{Phase 4: SRP Decommissioning}

Migration completion occurs when all customers and actions have successfully transitioned to SRM execution, validated through sustained periods (30+ days) of stable operation without SRP fallback. The SRP system transitions to read-only mode, the catalog synchronization job is deactivated, and finally the SRP infrastructure is decommissioned after a retention period ensuring all audit and compliance requirements are satisfied.

\section{Data Integrity and Consistency Mechanisms}

The coexistence architecture implements several mechanisms to maintain data integrity and consistency during the migration period.

\subsection{Transactional Boundaries and Idempotency}

Synchronization operations execute within database transaction contexts ensuring atomicity: either all records in a synchronization batch commit successfully, or none do. The synchronization logic is designed for idempotency, enabling safe re-execution in case of failures. Each sync function can be invoked multiple times with identical results, as record identification is based on deterministic primary keys from SRP.

\subsection{Referential Integrity and Audit Trail}

Foreign key relationships are validated before record creation to prevent orphaned references. The synchronization order ensures parent entities exist before children, while optional relationships are nulled when references cannot be resolved. While audit logging is disabled during bulk synchronization for performance, the synchronization jobs themselves are comprehensively logged through the monitoring infrastructure, with execution metrics recorded in New Relic.

\section{Operational Considerations and Migration Benefits}

The coexistence strategy introduces several operational challenges that require careful management during the migration period.

\subsection{Operational Challenges}

The primary challenges include: dual maintenance burden requiring catalog changes in SRP that propagate through synchronization; synchronization latency introducing up to 24-hour delays between SRP changes and SRM visibility; execution history fragmentation requiring staff training to query SRM exclusively; and performance considerations necessitating scheduled execution during low-traffic periods.

\subsection{Risk Mitigation and Benefits}

The coexistence and gradual migration strategy provides a pragmatic path to modernization while maintaining operational continuity. The dual-flow architecture enables incremental validation of the new system in production environments, building confidence through measured progression rather than high-risk wholesale replacement. The phased rollout provides operations teams with extended periods to develop familiarity with the new system architecture, monitoring tools, and troubleshooting procedures.

The coexistence strategy eliminates migration-related service disruptions by maintaining the legacy system as a fully functional fallback throughout the transition period. This approach ensures that service level agreements are maintained throughout the migration, avoiding the business impact of failed "big bang" migrations. The gradual migration provides real-world validation of architectural decisions under production loads, enabling refinement based on operational feedback before full deployment.

While the strategy introduces operational complexity through the need for synchronization and dual-system maintenance, these costs are justified by the substantial risk reduction and business continuity assurance provided by the gradual approach. The coexistence period represents a necessary investment in migration safety that protects both the company and its customers from the substantial risks inherent in large-scale system replacements.

\section{Production Deployment and Real-World Experience}

The transition from design and implementation to production deployment required a carefully orchestrated rollout strategy that validated the new architecture under real-world conditions while minimizing disruption to existing customer workflows.

\subsection{Initial Frontend Deployment}

The first SRM component deployed to production was the frontend application, representing a strategic decision to decouple user interface improvements from backend execution modernization. This deployment strategy enabled all service requests to flow through the new Vue.js 3 frontend, immediately benefiting from enhanced field validation logic, improved user experience, and the refined three-level configuration system, while maintaining backward compatibility with the legacy SRP execution pipeline.

This deployment implemented intelligent routing logic that directed service requests to either the new SRM backend or the legacy SRP system based on action configuration and customer migration status. For customers and actions not yet migrated to SRM, the frontend created service requests through SRP's existing REST API, ensuring transparent operation from the user perspective regardless of the underlying execution mechanism. This approach validated the frontend architecture independently, enabling rapid iteration on user interface issues without the complexity of simultaneous backend migration.

The decoupled deployment also provided immediate value to all customers through standardized field rendering, client-side validation improvements, and general performance optimizations, establishing user confidence in the modernization effort before introducing backend changes that would affect critical execution workflows.

\subsection{First Customer Migration: Irca Case Study}

Following successful frontend stabilization, the first complete backend migration targeted Irca, a strategic customer selection motivated by several converging factors. Irca had requested a new integration enabling programmatic service request creation from their internal business process automation platform, requiring API enhancements and customization work that necessitated close collaboration and frequent communication. This existing engagement provided natural opportunities for migration coordination, technical validation sessions, and rapid feedback loops that would be challenging to establish with customers not actively involved in development initiatives.

The Irca migration served as a comprehensive validation of the complete SRM architecture under production load, exercising the Ansible-based execution pipeline, the three-level catalog configuration hierarchy, asynchronous task processing through Celery, and integration with external systems including HDA ticket management and the Provision Prime provisioning platform. The migration proceeded through a staged rollout where subsets of Irca's action catalog were progressively transitioned to SRM execution, with each wave monitored for execution success rates, latency characteristics, and error patterns before expanding scope.

Technical challenges encountered during the Irca migration primarily centered on edge case error handling that had not been anticipated during development and testing phases. Several scenarios involving malformed input data, unexpected field value combinations, and race conditions in asynchronous task execution triggered HTTP 500 Internal Server Error responses that, while handled correctly from a system stability perspective, provided insufficient diagnostic information to end users and support staff. These issues were addressed through enhanced validation logic that detects edge cases earlier in the request lifecycle, custom exception handlers that translate technical errors into user-friendly messages, and comprehensive logging that enables rapid diagnosis of unexpected failure modes.
\subsection{Message Broker Architecture Refinement}

Production deployment experience revealed the critical importance of message queue reliability for the asynchronous task execution architecture. The initial implementation utilized Redis as the message broker for Celery task queues, a common configuration that offers simplicity and high performance \cite{redis_kafka_rabbitmq_comparison}. However, operational analysis identified message durability as a critical concern: Redis's in-memory architecture introduces risk of task loss in failure scenarios, particularly during broker restarts or infrastructure incidents.

Given that service request execution tasks represent business-critical operations with compliance and audit implications, message loss would constitute an unacceptable failure mode potentially resulting in service requests entering inconsistent states where database records indicate pending execution but no corresponding task exists in the queue. Such scenarios would require manual intervention to identify and remediate, introducing operational overhead and risking Service Level Agreement (SLA) violations.

The system was therefore migrated to RabbitMQ as the Celery message broker, a production-grade message queuing platform designed specifically for reliable message delivery in distributed systems. RabbitMQ provides several critical guarantees absent from Redis \cite{telecom4020018} \cite{redis_kafka_rabbitmq_comparison}: persistent message storage with configurable durability ensuring tasks survive broker restarts; publisher confirms enabling the application to verify successful message delivery before considering a task enqueued; consumer acknowledgments ensuring tasks are not removed from queues until worker processes confirm successful completion; and dead letter exchanges enabling systematic handling of tasks that cannot be processed \cite{telecom4020018}.

This architectural refinement reflects the maturity progression from development and testing environments, where Redis's simplicity and performance characteristics are appropriate, to production deployment where reliability and operational guarantees take precedence \cite{telecom4020018} \cite{redis_kafka_rabbitmq_comparison}. The migration to RabbitMQ introduced minimal application-level changes due to Celery's broker abstraction layer, primarily requiring configuration updates and deployment of RabbitMQ infrastructure within the Kubernetes cluster.

\subsubsection{RabbitMQ}
The production deployment ultimately adopted RabbitMQ as the message broker, driven by its superior reliability guarantees and enterprise-grade resilience features. RabbitMQ implements the Advanced Message Queuing Protocol (AMQP) \cite{amqp}, which provides strong delivery semantics through acknowledgment-based message processing: tasks are not removed from queues until workers explicitly confirm successful completion, ensuring zero message loss even during worker crashes or network failures \cite{amqp}. This transactional model, combined with RabbitMQ's durable queues and persistent message storage, guarantees that critical service orchestration tasks survive broker restarts and system failures, a crucial requirement for production environments where service provisioning must be reliably tracked and completed.

RabbitMQ's sophisticated routing capabilities through exchanges, bindings, and multiple queue types (classic, quorum, streams) enable complex message distribution patterns that scale effectively across distributed worker pools \cite{videla2012rabbitmq}. The broker's clustering and federation features support high-availability deployments with automatic failover, while its flow control mechanisms prevent memory exhaustion by applying backpressure when consumers cannot keep pace with producers. Additionally, RabbitMQ's mature monitoring ecosystem, including native management UI and Prometheus integration, provides comprehensive visibility into queue depths, throughput rates, and system health—essential operational capabilities for managing asynchronous task processing at scale \cite{videla2012rabbitmq}. These enterprise-focused features, though unnecessary during initial development, proved indispensable for production reliability and operational observability.

\subsection{Performance Optimization and Iterative Refinement}

Production telemetry from the initial deployments revealed performance bottlenecks in the action detail retrieval endpoint used by MyElmec to render service request creation forms. The original implementation pre-computed all metadata required for form rendering, including external field configurations for HDA ticket classification, before returning the API response. For actions with complex classification requirements involving multiple external API calls to retrieve categorization options, this synchronous approach introduced latencies exceeding 3-5 seconds, degrading perceived responsiveness and user experience.

The optimization strategy restructured the workflow to implement progressive rendering where the backend returns the core action definition and field configurations immediately, enabling the frontend to render the primary form interface without delay. External HDA classification fields, which are not required for initial form display, are fetched asynchronously by the frontend after the form renders, populating classification dropdowns in the background while users begin completing primary fields. This architectural change reduced median response times from the action detail endpoint to under 500ms, substantially improving perceived performance.

Security considerations required that backend validation of HDA classification fields remain comprehensive despite the asynchronous frontend fetching strategy. The backend continues to validate classification field values during service request submission, ensuring that malicious or malformed classification data cannot bypass validation through client-side manipulation. This defense-in-depth approach maintains security guarantees while optimizing the critical path of form rendering for legitimate users.

\subsection{Current Migration Status and Future Timeline}

As of January 2026, the SRM system operates in production serving IRCA and several additional customers migrated through the progressive expansion phase described in Section 4.5. These customers collectively represent diverse action catalog configurations, execution patterns, and integration requirements, providing comprehensive validation of the architecture's flexibility and scalability. The legacy SRP system continues to serve the majority of the customer base, with catalog synchronization and execution history merging maintaining operational continuity during the coexistence period.

Customer feedback collected through support channels, direct communication with pilot customers, and implicit signals from system usage patterns has been predominantly positive. While much of the innovation introduced by SRM operates transparently at the infrastructure and operational level rather than through visible user-facing features, customers have noted improvements in service request creation workflow smoothness, reduction in validation errors requiring form resubmission, and faster response times during form interaction.

The migration roadmap targets complete SRP decommissioning by the end of Q2 2026, a timeline driven by progressive customer migration velocity and validation of operational stability at increasing scale. Each migration wave incorporates learnings from previous deployments, refining the migration playbook and reducing per-customer migration effort through automation and standardization. The aggressive but achievable timeline reflects confidence in the architecture's production readiness established through successful operation of migrated customers without significant incidents or rollback requirements.

Operational teams have developed proficiency with the new monitoring infrastructure, troubleshooting workflows, and catalog management interfaces through extended exposure during the coexistence period. This operational readiness, combined with demonstrated technical stability and positive customer reception, positions the organization to execute the remaining migrations and achieve full transition to the modern SRM architecture within the planned timeframe.